\section{The volume problem and its normalization}
\label{sec:conventions}

\paragraph{What is being measured?}
A moduli space is a space whose points are geometric objects, with
objects related by the prescribed equivalence regarded as the same point.
Here one point is an oriented hyperbolic surface, considered up to
orientation-preserving isometry that preserves the labels of its ends.
The genus $g$ counts handles; $n$ counts boundary components or cusps.
A geodesic boundary is a closed curve that is straight for the intrinsic
metric. A cusp is a noncompact end with finite area, modeled by a
horoball modulo a translation. ``Unparametrized'' means that a boundary
has a label and a length but no chosen starting point.

The volume of this \emph{space of surfaces} is different from the area
of one surface. For fixed boundary lengths, the latter is already
$2\pi(2g-2+n)$ by Gauss--Bonnet. The volume problem asks how much
room there is to vary the shape. The Weil--Petersson (WP) symplectic
form supplies the measure for that question. A symplectic form is a
closed, nondegenerate alternating two-form; on a space of real
dimension $2d$, its top exterior power $\omega^d/d!$ is the associated
volume form. The integer $d$ below is half the real dimension, and
$\lengths$ abbreviates the ordered list $(L_1,\ldots,L_n)$.

Let $\M_{g,n}(L_1,\ldots,L_n)$ denote the moduli space of oriented hyperbolic
surfaces of genus $g$ with $n$ labeled, unparametrized geodesic boundary
components of prescribed lengths. A zero length denotes a cusp. Throughout,
the curvature is $-1$, the parameters are fixed before taking a volume, and
$2g-2+n>0$. Put
\[
 d=3g-3+n,\qquad
 V_{g,n}(\lengths)=\int_{\M_{g,n}(\lengths)}
                  \frac{\omega_{\WP}^{d}}{d!}.
\]
A useful way to see the dimension is to cut the surface into
\emph{pairs of pants}, meaning spheres with three holes. The
$3g-3+n$ interior cutting curves each contribute a length and a
relative gluing displacement, or twist. These are the
Fenchel--Nielsen coordinates $(\ell_j,\tau_j)$ on the marked
space: a marking identifies the surface topologically with a fixed
reference surface. Forgetting that marking takes the quotient by
mapping classes, that is, orientation-preserving homeomorphisms up
to isotopy. The length--twist coordinates are local coordinates on
moduli, with possible finite symmetries. They do not make the whole
moduli space a rectangular region.

For a pair of pants, $(g,n)=(0,3)$ and $d=0$: the three boundary
lengths determine its metric, so there is no shape parameter left.
For a one-holed torus, $(g,n)=(1,1)$ and $d=1$: cutting one curve
leaves a pair of pants with boundary lengths $(L,\ell,\ell)$,
and gluing the two equal boundaries requires the twist $\tau$.

We use the Fenchel--Nielsen convention
\begin{equation}
 \omega_{\WP}=\sum_{j=1}^{d}\dd\ell_j\wedge\dd\tau_j,
 \qquad \text{one full Dehn twist: }\tau_j\longmapsto\tau_j+\ell_j.
 \label{eq:fn}
\end{equation}
This fixes the scale of the form independently of the convention for
orbifold integration. For comparison with complex coordinates $t_j$, recall that an
infinitesimal change of complex structure can be represented locally
by a \emph{Beltrami differential} $\mu(w)\,\dd\bar w/\dd w$.
The harmonic representative is the canonical representative selected
by the hyperbolic metric; its coefficient $\mu(w)$ is what enters
the following integral. Let
$\mu_j$ be the harmonic Beltrami representative of $\partial/\partial t_j$
and let $H_{j\bar k}=\int_X\mu_j\overline{\mu_k}\,\dd A$ use the curvature
$-1$ area element. Our positive complex convention is
\[
 \omega_{\WP}=i\sum_{j,k}H_{j\bar k}\,\dd t_j\wedge\dd\bar t_k.
\]
The factor is $i$, rather than $i/2$ with the same raw coefficients.
For closed surfaces the conversion is explicit in
\cite[Section 6.1, equation (38)]{KawazumiWP}; that paper writes
$\dd\tau\wedge\dd\ell$, so its twist direction is opposite to ours.
The conical extension used in Section~\ref{sec:flatlimit} retains this
scale; its geometric hypotheses are stated there.

Volumes are orbifold, or stack, volumes: generic ineffective stabilizers are
included. In particular,
\begin{equation}
 V_{0,3}=1,\qquad
 V_{0,4}(\lengths)=2\pi^2+\frac12\sum_{i=1}^4L_i^2,
 \qquad
 V_{1,1}(L)=\frac{L^2+4\pi^2}{48}.
 \label{eq:basevolumes}
\end{equation}
The effective torus volume is twice the last expression. This convention
matters both in unfolding and in the torus--sphere correspondence; see
Do--Norbury~\cite{DoNorbury} and Section~\ref{sec:torussphere}.
More generally, a generic ineffective stabilizer of order $h$ divides the
effective volume by $h$ without rescaling the local two-form. Rescaling
$\omega_{\WP}$ by a constant $c$ instead multiplies a $2d$-dimensional
volume by $c^d$. These are separate normalization operations.

The orbifold convention can be understood in a finite quotient chart.
If a finite symmetry group of order $h$ acts on a parameter region,
its orbifold integral is the integral upstairs divided by $h$.
This includes elements that act trivially on the parameters, called
\emph{ineffective} symmetries. On the one-holed torus the generic
elliptic involution contributes such a factor of two. Thus the
stack volume at $L=0$ is $\pi^2/12$, whereas the effective quotient
volume is $\pi^2/6$. This changes the counting of objects with
symmetries; it does not replace the local form $\dd\ell\wedge\dd\tau$
by half that form.

The original geometric proposal has three separate ingredients:
\begin{enumerate}
 \item realize an intrinsic surface by a canonical convex polyhedron;
 \item describe the allowed polyhedra by cells, ideally linear polytopes;
 \item identify the pullback of $\omega_{\WP}^{d}/d!$ on each cell.
\end{enumerate}
The realization and the cell decomposition identify the integration domain.
They do not, by themselves, identify its measure. In the ideal genus-zero
case the measure comparison is particularly favorable; in the hyperideal
case a useful local density is still available, but it is generally no longer
constant in the natural angle variables.

There are also two different cross-ratios in the four-point problem.
The cross-ratio is the complex parameter obtained by taking three
ordered points on the Riemann sphere to $0,1,\infty$ by a M\"obius
transformation and recording the image of the fourth, with the
ordering fixed. Four
ideal vertices on $\partial\HH^3$ have an \emph{ambient} cross-ratio. The
intrinsic metric on the polyhedral boundary has a conformal uniformization,
and its four punctures have a \emph{conformal} cross-ratio. The realization
theorem relates the two parameter spaces, but these two numerical coordinates
are not generally equal. Keeping them distinct is essential to the first
explicit integral below.

The computational companion in
\texttt{experiments/normalizations/README.md} records exact finite checks
of the determinants, factors, and low-dimensional intersection formulas
used in Sections~\ref{sec:conventions}--\ref{sec:flatlimit} and
Appendix~\ref{app:fricke}; its machine-readable output is
\path{artifacts/normalizations/result.json}. These fixtures check
algebraic consequences and conventions. They do not prove the cited
realization, measure-comparison, or convergence theorems.

